\Needspace{20\baselineskip}
\section{吞吐率、屋顶线与存储流量}
\label{l14:sec:roofline}

\subsection{理论峰值与资源利用指标}

硬件的理论吞吐能力常被称为峰值上限（Speed of Light，SoL）。内核报告中的计算和存储汇总指标，用相对于对应资源峰值的比例描述当前活动强度。由于一个汇总项包含多个组成资源，应先检查它由哪些计数器形成，再判断接近饱和的是浮点流水线、加载／存储流水线、缓存入口还是设备内存接口。

设某个汇总指标的组成资源利用比例为 $u_1,\ldots,u_q$，相关吞吐指标采用最大组成比例，即 $u_{\mathrm{group}}=\max_j u_j$。例如，加载／存储通路达到约 $98\%$，而浮点流水线只有较低活动度，汇总项仍可能接近 $98\%$。它描述的是最忙组成资源的状态，不能直接用来计算全芯片的总体浮点效率。

这与式\eqref{l14:eq:performance}的效率 $\eta=P/P_{\mathrm{peak}}$ 有不同的用途。$\eta$ 由选定的浮点工作量和耗时计算，回答目标运算相对于对应算力峰值达到了多少；资源汇总及其分解则帮助判断什么硬件正在承受最大压力。若两者相差很大，通常应继续检查计数口径、精度、时钟，以及非算术指令的开销。

\subsection{算术强度与屋顶线模型}

\term{算术强度}{arithmetic intensity}表示每传输一个字节所完成的浮点运算量。设某内核工作量为 $F$，在指定存储层级 $\ell$ 实际读取和写入的总字节数为 $Q_\ell$，则 $I_\ell=F/Q_\ell$，单位为 FLOP/byte。存储层级必须明确：设备内存的读写量、L2 与 L1 之间的流量、共享内存访问量，各自对应不同的强度和带宽上限。

先考虑一个只有计算峰值 $P_{\mathrm{peak}}$ 与指定层级带宽峰值 $B_\ell$ 的模型。完成全部计算至少需要 $F/P_{\mathrm{peak}}$ 时间，完成全部数据传输至少需要 $Q_\ell/B_\ell$ 时间。即使二者理想重叠，总时间也不能小于其中的较大者。因此
\begin{equation}
 t\geq\max\left(\frac{F}{P_{\mathrm{peak}}},
                 \frac{Q_\ell}{B_\ell}\right)
 \quad\Longrightarrow\quad
 P=\frac{F}{t}\leq\min\left(P_{\mathrm{peak}},B_\ell I_\ell\right).
 \label{l14:eq:roofline}
\end{equation}
这个上界称为\term{屋顶线模型}{roofline model}。在双对数坐标下，$B_\ell I_\ell$ 是斜线，$P_{\mathrm{peak}}$ 是水平线；二者相交处的算术强度为 $I_\ell^*=P_{\mathrm{peak}}/B_\ell$，也称转折点。

\begin{figure}[htbp]
\centering
\begin{tikzpicture}[x=1cm,y=1cm,>=Latex,font=\small]
  \draw[->] (0,0)--(11.7,0) node[right] {$I_\ell$};
  \draw[->] (0,0)--(0,4.7) node[above] {$P$};
  \draw[line width=1.3pt,color=UIUCOrange] (0.5,0.5)--(5.7,3.8);
  \draw[line width=1.3pt,color=UIUCBlue] (5.7,3.8)--(11,3.8);
  \draw[color=black!35] (5.7,0)--(5.7,3.8);
  \fill (5.7,3.8) circle(1.8pt);
  \node[below] at (5.7,0) {$I_\ell^*$};
  \node[above] at (8.7,3.8) {计算峰值 $P_{\mathrm{peak}}$};
  \node[above,rotate=32] at (3.0,2.28) {带宽上界 $B_\ell I_\ell$};
  \fill[color=UIUCBlue] (3.3,1.35) circle(2.8pt);
  \node[right] at (3.5,1.35) {实际性能};
  \node at (2.5,0.45) {带宽上界更低};
  \node at (8.7,0.45) {计算上界更低};
\end{tikzpicture}
\caption{屋顶线模型的双对数示意。上界由算力和指定层级带宽中更严格的一项确定，实际内核还可能受到延迟、调度或指令通路限制。}
\label{l14:fig:roofline-model}
\end{figure}

当 $I_\ell<I_\ell^*$ 时，带宽项更小；当 $I_\ell>I_\ell^*$ 时，计算项更小。在内核具有足够并行度、能够充分驱动相关资源的条件下，二者分别对应\term{带宽受限}{memory-bound}和\term{计算受限}{compute-bound}的典型区域。若测量点离两条上界都很远，则需要进一步查找延迟、依赖、指令发射或其他片上通路的限制；仅比较横坐标与转折点不足以给出完整诊断。

例如，设某 GPU 对当前精度的峰值为 $20$ TFLOP/s，对应层级带宽为 $1$ TB/s，则转折点为 $20$ FLOP/byte。算术强度为 $4$ 的内核，其模型上界为 $4$ TFLOP/s；若实测只有 $0.5$ TFLOP/s，就仅达到该上界的 $12.5\%$，还需检查为什么带宽没有被有效使用。若强度提高到 $80$，带宽项达到 $80$ TFLOP/s，模型上界则由 $20$ TFLOP/s 的计算峰值决定。这些数字用于说明计算步骤，并非某一实验 GPU 的规格。

不同 GPU 的算力和带宽比例不同，同一个内核的横坐标也可能因缓存行为改变。因此，换用计算峰值更高的 GPU 后，内核可能更接近带宽限制；反过来，带宽提高较多的设备也可能让计算上界成为更严格的约束。使用多个精度的屋顶线时，还需要选择与实际执行精度相符的水平峰值线。

\subsection{矩阵乘法中不同层级的复用}

算术强度的推导必须说明数据流量假设。对单精度 $C=AB$，若 $A$ 和 $B$ 各从设备内存读入一次，$C$ 写回一次，则理想流量为 $Q_{\mathrm{once}}=4(MK+KN+MN)$ 字节，理想强度为
\begin{equation}
 I_{\mathrm{once}}=\frac{2MNK}{4(MK+KN+MN)}.
 \label{l14:eq:gemm-ai}
\end{equation}
对边长为 $n$ 的方阵，这给出 $I_{\mathrm{once}}=n/6$。这个结果描述一次读入、一次写回的理想流量；真实分块中，不同线程块可能重复读取输入子块，能否由缓存承接这些重复访问取决于实际执行。

为了比较源代码层面的加载需求，可暂时忽略缓存，并令 $n$ 能被分块宽度 $T$ 整除。简单内核有约 $2n^3$ 次单精度输入加载和 $n^2$ 次输出写入，因此逻辑字节数约为 $8n^3+4n^2$，强度在 $n$ 较大时趋近 $1/4$。共享分块将输入加载减少到约 $2n^3/T$ 个元素，逻辑字节数变为 $8n^3/T+4n^2$，相应强度趋近 $T/4$。例如 $T=16$ 时，这个逻辑加载模型给出约 $4$ FLOP/byte。

这里的 $1/4$、$T/4$ 与 $n/6$ 采用了不同的数据流量假设。前两者比较内核发出的逻辑全局访问需求，后者是理想的一次读写边界。若使用设备内存屋顶线，就应以实测设备内存字节数计算 $I_{\mathrm{DRAM}}$；缓存命中的重复加载不会每次都计入这部分字节。因此，简单内核在实测设备内存屋顶线中的强度完全可能高于共享内核，同时运行得更慢。

在带宽上界主导且其他资源仍有余量时，提高每份已加载数据上的有效计算量，例如复用数据或融合需要相同输入的运算，可以提高强度。新增计算仍要满足寄存器、同步和指令通路的限制；人为加入无用运算，只会改变统计工作量，不能构成有意义的应用加速。对已经接近计算峰值的内核，则应考虑减少必要运算、改善指令组织，或在精度要求允许时使用更适合的计算路径。低精度路径的峰值提升幅度取决于具体架构和指令，不采用固定倍数作为通用结论。
